% \documentclass[letterpaper,12pt]{article}
\documentclass[letterpaper,12pt]{book}
\usepackage{amsmath}
\usepackage{amssymb}
\usepackage{amsfonts}
\usepackage{array}
\usepackage{longtable}

\usepackage{framed}
% \usepackage{subfig}
\usepackage{multirow}
\usepackage{hyperref}

%% graphics and images
\usepackage{lscape}
\usepackage{graphicx}
\usepackage{graphics}
\usepackage{rotating}
\usepackage{caption}
\usepackage{subcaption}
% \usepackage{multirow}

\usepackage{setspace}
\usepackage{titlesec}

%% Algorithm package
% \usepackage{algorithm}
\usepackage[chapter]{algorithm}
\usepackage[noend]{algpseudocode}

%% Algorithm custom command
\makeatletter
\def\BState{\State\hskip-\ALG@thistlm}
\makeatother

%% verbatim samepage
\usepackage{fancyvrb}

\usepackage{xspace}

%% \usepackage{natbib}
%% \bibliographystyle{abbrvnat}
%% \bibpunct{(}{)}{;}{a}{,}{,}
\usepackage[utf8]{inputenc}

\usepackage[spanish]{babel}

\setlength\topmargin{-0.375in}
\setlength\headheight{0in}
\setlength\headsep{0.375in}
\setlength\textheight{9.5in}
\setlength\textwidth{7.0in}
\setlength\oddsidemargin{-0.25in}
\setlength\evensidemargin{-0.25in}
\setlength\parindent{0.25in}
\setlength\parskip{0.25in}

%% \usepackage{amssymb}
%% using square for first level listing instead circles
\renewcommand{\labelitemi}{\tiny$\blacksquare$}

%% \usepackage{cite}
\usepackage[backend=bibtex, bibencoding=utf8]{biblatex}

\usepackage[letterpaper, portrait, margin=1.25in]{geometry}


\providecommand{\lessgtr}{\stackrel{<}{>}} %% could be provided natively with different font encoding:
\IfFileExists{pslatex.sty}{

  \usepackage{pslatex}
}{}

%% Custom commands
\newcommand{\eng}[1]{\textit{#1}\xspace}


\newcolumntype{M}[1]{>{\centering\arraybackslash}m{#1}}
\newcolumntype{N}{@{}m{0pt}@{}}




\DefineBibliographyStrings{english}{
  references = {Bibliografía},
}
\bibliography{references.bib}


\titleformat{\chapter}[display]
{\normalfont\huge\bfseries}{Capítulo \ \thechapter}{5pt}{\Huge}   
\titlespacing*{\chapter}{0pt}{-50pt}{40pt}
\titlespacing*{\section}{0pt}{0pt}{0pt}



\begin{document}

\begin{titlepage}
\newgeometry{top=3cm}
\begin{centering}
{
\textbf{\Large UNIVERSIDAD TÉCNICA FEDERICO SANTA MARÍA}\\[0em]
\textbf{\small DEPARTAMENTO DE INFORMÁTICA}\\[0em]
\textbf{\small SANTIAGO -- CHILE}\\[2.4em]
}
\includegraphics[scale=0.295]{content/figures/logo-utfsm} \\[4em]
\begin{spacing}{1.8}
\textbf{\Large \uppercase{``Título de la memoria''}}\\[3.70em]
\end{spacing}
\textbf{\large \uppercase{Pablo Antonio Ortega Mesa}}\\[0.54em]
\textbf{\small MEMORIA DE TITULACIÓN PARA OPTAR AL TÍTULO DE INGENIERO CIVIL INFORMÁTICA} \\[1em]
\begin{tabular}{ccc} \small\bf
\small\bf PROFESOR GUÍA:          & & \small\bf \uppercase{Marcelo Mendoza} \\
\small\bf PROFESOR CORREFERENTE:  & & \small\bf \uppercase{Profesor Correferente} \\
\end{tabular} \\[2.5em]
\textbf{\small Mayo -- 2015} \\
\end{centering}
\end{titlepage}


\setstretch{1.5}

\section*{Resumen}
Lorem ipsum dolor sit amet, consectetur adipisicing elit, sed do eiusmod
tempor incididunt ut labore et dolore magna aliqua. Ut enim ad minim veniam,
quis nostrud exercitation ullamco laboris nisi ut aliquip ex ea commodo
consequat. Duis aute irure dolor in reprehenderit in voluptate velit esse
cillum dolore eu fugiat nulla pariatur. Excepteur sint occaecat cupidatat non
proident, sunt in culpa qui officia deserunt mollit anim id est laborum.
\clearpage

\section*{Abstract}
Lorem ipsum dolor sit amet, consectetur adipisicing elit, sed do eiusmod
tempor incididunt ut labore et dolore magna aliqua. Ut enim ad minim veniam,
quis nostrud exercitation ullamco laboris nisi ut aliquip ex ea commodo
consequat. Duis aute irure dolor in reprehenderit in voluptate velit esse
cillum dolore eu fugiat nulla pariatur. Excepteur sint occaecat cupidatat non
proident, sunt in culpa qui officia deserunt mollit anim id est laborum.
\clearpage

\section*{Agradecimientos}
Lorem ipsum dolor sit amet, consectetur adipisicing elit, sed do eiusmod
tempor incididunt ut labore et dolore magna aliqua. Ut enim ad minim veniam,
quis nostrud exercitation ullamco laboris nisi ut aliquip ex ea commodo
consequat. Duis aute irure dolor in reprehenderit in voluptate velit esse
cillum dolore eu fugiat nulla pariatur. Excepteur sint occaecat cupidatat non
proident, sunt in culpa qui officia deserunt mollit anim id est laborum.
\clearpage

\let\cleardoublepage\clearpage

%% comment it out if you don't like it and/or use one of the lines below:
%% \setstretch{2}
%\setstretch{1.5}

\tableofcontents
\clearpage

% \listoffigures
% \clearpage

% \listoftables
% \clearpage

% \newpage
\chapter{Introducción}
\label{chap:intro}
En esta sección se propone escribir sobre la importancia del descubrimiento de conocimiento por parte de la comunidad científica. El principal problema para esta tarea es la búsqueda dentro de millones de documentos científicos sin una clasificación correcta con el objetivo de obtener la literatura relacionada con las investigaciones en curso. 

Por otro lado se planteará la importancia de disponer de una correcta visualización para la clasificación de documentos y las utilidades de esta.

Se hará referencia las técnicas FCA y LDA. Además se mencionará el uso que se le da en este trabajo.

Finalmente, se planteará la propuesta de este trabajo y se plantearán los objetivos.

\clearpage

\chapter{Marco Teórico}
\label{chap:mteorico}
En este capítulo se hablará sobre el marco teórico de las técnicas utilizadas: FCA y LDA. Además de hacer un breve análisis de los trabajos relacionados a cada una de estas.

\section{Formal Concept Analysis}
\label{sec:fca}

En esta sección se planteará todo el framework matemático que constituye FCA. Luego, se realizará un breve análisi de la literatura asociada. (NOTA: utilizar la referencia utilizada por Víctor Codocedo en presentación, ya que incluye los principales ``milestones'' de la técnica).

\section{Topic Modeling}
\label{sec:topicmodeling}

En esta sección se describirá en que consiste Topic Modeling, sus objetivos y como nace. Haciendo una breve pasada por \emph{PSLI}.

\subsection{Modelo LDA}
\label{subsec:ldamodel}
Se explicará en detalle en que consiste el modelo LDA, presentando ejemplos para la mejor comprensión del lector.

\subsection{Inferencia Estadística}
\label{subsec:infest}
Para continuar explicando \emph{Modelo LDA}

\subsection{Gibbs Sampling}
\label{subsec:gibbssampling}
Se explicará la técnica \emph{Gibbs Sampling} ya que es la utilizada por los algoritmos de este trabajo.

\section{DBLP}
\label{sec:dblp}
DBLP es una gran colección de trabajos científicos muy importante para la comunidad. Se explicará la importancia, los desafíos que arroja al ser una ``\emph{Gran Colección}'' y el porqué se escogió esta base de datos.
\clearpage

\chapter{Descripción del Problema}
\label{chap:problem}
Lorem ipsum dolor sit amet, consectetur adipisicing elit, sed do eiusmod
tempor incididunt ut labore et dolore magna aliqua. Ut enim ad minim veniam,
quis nostrud exercitation ullamco laboris nisi ut aliquip ex ea commodo
consequat. Duis aute irure dolor in reprehenderit in voluptate velit esse
cillum dolore eu fugiat nulla pariatur. Excepteur sint occaecat cupidatat non
proident, sunt in culpa qui officia deserunt mollit anim id est laborum.

\section{Base de Datos Científicas}
\label{sec:sciencedb}
Lorem ipsum dolor sit amet, consectetur adipisicing elit, sed do eiusmod
tempor incididunt ut labore et dolore magna aliqua. Ut enim ad minim veniam,
quis nostrud exercitation ullamco laboris nisi ut aliquip ex ea commodo
consequat. Duis aute irure dolor in reprehenderit in voluptate velit esse
cillum dolore eu fugiat nulla pariatur. Excepteur sint occaecat cupidatat non
proident, sunt in culpa qui officia deserunt mollit anim id est laborum.

\section{Dimensionalidad}
\label{sec:dimensionality}
Lorem ipsum dolor sit amet, consectetur adipisicing elit, sed do eiusmod
tempor incididunt ut labore et dolore magna aliqua. Ut enim ad minim veniam,
quis nostrud exercitation ullamco laboris nisi ut aliquip ex ea commodo
consequat. Duis aute irure dolor in reprehenderit in voluptate velit esse
cillum dolore eu fugiat nulla pariatur. Excepteur sint occaecat cupidatat non
proident, sunt in culpa qui officia deserunt mollit anim id est laborum.

\section{Visualización de Lattices}
\label{sec:latticevisualization}
Lorem ipsum dolor sit amet, consectetur adipisicing elit, sed do eiusmod
tempor incididunt ut labore et dolore magna aliqua. Ut enim ad minim veniam,
quis nostrud exercitation ullamco laboris nisi ut aliquip ex ea commodo
consequat. Duis aute irure dolor in reprehenderit in voluptate velit esse
cillum dolore eu fugiat nulla pariatur. Excepteur sint occaecat cupidatat non
proident, sunt in culpa qui officia deserunt mollit anim id est laborum.
\clearpage

\chapter{Solución Propuesta}
\label{chap:solution}
Lorem ipsum dolor sit amet, consectetur adipisicing elit, sed do eiusmod
tempor incididunt ut labore et dolore magna aliqua. Ut enim ad minim veniam,
quis nostrud exercitation ullamco laboris nisi ut aliquip ex ea commodo
consequat. Duis aute irure dolor in reprehenderit in voluptate velit esse
cillum dolore eu fugiat nulla pariatur. Excepteur sint occaecat cupidatat non
proident, sunt in culpa qui officia deserunt mollit anim id est laborum.

\section{Topic Modeling}
\label{sec:soltopicmodeling}
Lorem ipsum dolor sit amet, consectetur adipisicing elit, sed do eiusmod
tempor incididunt ut labore et dolore magna aliqua. Ut enim ad minim veniam,
quis nostrud exercitation ullamco laboris nisi ut aliquip ex ea commodo
consequat. Duis aute irure dolor in reprehenderit in voluptate velit esse
cillum dolore eu fugiat nulla pariatur. Excepteur sint occaecat cupidatat non
proident, sunt in culpa qui officia deserunt mollit anim id est laborum.

\section{FCA}
\label{sec:solfca}
Lorem ipsum dolor sit amet, consectetur adipisicing elit, sed do eiusmod
tempor incididunt ut labore et dolore magna aliqua. Ut enim ad minim veniam,
quis nostrud exercitation ullamco laboris nisi ut aliquip ex ea commodo
consequat. Duis aute irure dolor in reprehenderit in voluptate velit esse
cillum dolore eu fugiat nulla pariatur. Excepteur sint occaecat cupidatat non
proident, sunt in culpa qui officia deserunt mollit anim id est laborum.

\section{Visualización}
\label{sec:solvisualization}
Lorem ipsum dolor sit amet, consectetur adipisicing elit, sed do eiusmod
tempor incididunt ut labore et dolore magna aliqua. Ut enim ad minim veniam,
quis nostrud exercitation ullamco laboris nisi ut aliquip ex ea commodo
consequat. Duis aute irure dolor in reprehenderit in voluptate velit esse
cillum dolore eu fugiat nulla pariatur. Excepteur sint occaecat cupidatat non
proident, sunt in culpa qui officia deserunt mollit anim id est laborum.
\clearpage

\chapter{Resultados}
\label{chap:results}
Lorem ipsum dolor sit amet, consectetur adipisicing elit, sed do eiusmod
tempor incididunt ut labore et dolore magna aliqua. Ut enim ad minim veniam,
quis nostrud exercitation ullamco laboris nisi ut aliquip ex ea commodo
consequat. Duis aute irure dolor in reprehenderit in voluptate velit esse
cillum dolore eu fugiat nulla pariatur. Excepteur sint occaecat cupidatat non
proident, sunt in culpa qui officia deserunt mollit anim id est laborum.




\clearpage

\chapter{Trabajo futuro}
\label{chap:futurework}
\section*{Trabajo futuro 1}

Lorem ipsum dolor sit amet, consectetur adipisicing elit, sed do eiusmod
tempor incididunt ut labore et dolore magna aliqua. Ut enim ad minim veniam,
quis nostrud exercitation ullamco laboris nisi ut aliquip ex ea commodo
consequat. Duis aute irure dolor in reprehenderit in voluptate velit esse
cillum dolore eu fugiat nulla pariatur. Excepteur sint occaecat cupidatat non
proident, sunt in culpa qui officia deserunt mollit anim id est laborum.


\section*{Trabajo futuro 2}

Lorem ipsum dolor sit amet, consectetur adipisicing elit, sed do eiusmod
tempor incididunt ut labore et dolore magna aliqua. Ut enim ad minim veniam,
quis nostrud exercitation ullamco laboris nisi ut aliquip ex ea commodo
consequat. Duis aute irure dolor in reprehenderit in voluptate velit esse
cillum dolore eu fugiat nulla pariatur. Excepteur sint occaecat cupidatat non
proident, sunt in culpa qui officia deserunt mollit anim id est laborum.
\clearpage

\chapter{Conclusiones}
\label{chap:concl}
Lorem ipsum dolor sit amet, consectetur adipisicing elit, sed do eiusmod
tempor incididunt ut labore et dolore magna aliqua. Ut enim ad minim veniam,
quis nostrud exercitation ullamco laboris nisi ut aliquip ex ea commodo
consequat. Duis aute irure dolor in reprehenderit in voluptate velit esse
cillum dolore eu fugiat nulla pariatur. Excepteur sint occaecat cupidatat non
proident, sunt in culpa qui officia deserunt mollit anim id est laborum.

\section*{Conclusión 1}
Lorem ipsum dolor sit amet, consectetur adipisicing elit, sed do eiusmod
tempor incididunt ut labore et dolore magna aliqua. Ut enim ad minim veniam,
quis nostrud exercitation ullamco laboris nisi ut aliquip ex ea commodo
consequat. Duis aute irure dolor in reprehenderit in voluptate velit esse
cillum dolore eu fugiat nulla pariatur. Excepteur sint occaecat cupidatat non
proident, sunt in culpa qui officia deserunt mollit anim id est laborum.

\section*{Conclusión 2}
Lorem ipsum dolor sit amet, consectetur adipisicing elit, sed do eiusmod
tempor incididunt ut labore et dolore magna aliqua. Ut enim ad minim veniam,
quis nostrud exercitation ullamco laboris nisi ut aliquip ex ea commodo
consequat. Duis aute irure dolor in reprehenderit in voluptate velit esse
cillum dolore eu fugiat nulla pariatur. Excepteur sint occaecat cupidatat non
proident, sunt in culpa qui officia deserunt mollit anim id est laborum.
\newpage

% \bibliographystyle{plain}
\setstretch{1}
\nocite{*}
\printbibliography

\end{document}